\chapter{مفعول دوبلر}\label{ch:g12:doppler-effect}

تقترب سيارة إسعاف وصفّارتها تصرخ، ثم تمرّ --- وفي تلك اللحظة تنزلق النغمة
إلى أسفل انزلاقًا مسموعًا، مع أن السائق لم يمسّ شيئًا. فالصفّارة أمينة؛ أما \omterm{def:g10:signals-and-waves:wave}{الموجة}
فلا: إذ تعصرها الحركة أمامها وتمدّها خلفها. ويقيس هذا الفصل العصر،
ثم يركب الفكرة نفسها من الصفّارات إلى مسدسات \omterm{met:g10:signals-and-waves:echo}{الرادار}، إلى الدم في شريان، وخارجًا إلى
كوكب يفضحه تمايل في ضوء نجم.

\section{المفعول}

\begin{definition}[مفعول دوبلر]\label{def:g12:doppler-effect:doppler}
\emph{مفعول دوبلر}\index{مفعول دوبلر} هو تغيّر \omterm{def:g10:signals-and-waves:frequency}{التواتر} المستقبَل
\omterm{def:g10:signals-and-waves:wave}{لموجة} حين يتحرك المنبع والمستقبِل أحدهما بالنسبة إلى الآخر: \omterm{def:g10:signals-and-waves:frequency}{فالتواتر} المستقبَل
$f'$ أعلى من \omterm{def:g10:signals-and-waves:frequency}{التواتر} المصدَر $f$ ما داما يقتربان، وأدنى ما داما
يتباعدان، ولا يساوي $f$ إلا ما دامت المسافة بينهما تكفّ للحظة عن التغير.
\end{definition}

\begin{remark}[الحدّة لا العلوّ]\label{rem:g12:doppler-effect:pitch}
يتغير أمران عند مرور سيارة الإسعاف. أما \emph{\omterm{def:g12:sound-acoustics:pitch-loudness}{العلوّ}} فيتضخم ويخبو لأن
\omterm{def:g10:signals-and-waves:wave}{الموجة} تنتشر مع المسافة (\cref{ch:g12:sound-acoustics})؛ أما \emph{\omterm{def:g12:sound-acoustics:pitch-loudness}{الحدّة}} فتنزاح
بسبب سرعة الاقتراب أو الابتعاد، وتهبط فجأةً عند المرور.
\omterm{def:g12:sound-acoustics:pitch-loudness}{فالعلوّ} يقول ``كم البعد''؛ \omterm{def:g12:sound-acoustics:pitch-loudness}{والحدّة} تقول ``كم السرعة''.
\end{remark}

\section{منبع متحرك: جبهات موجة متزاحمة}

الصورة المفتاح: متى غادرت قمةٌ المنبعَ، ملكها الوسط. فتتوسع كل قمة
دائرةً بسرعة \omterm{def:g10:signals-and-waves:wave}{الموجة} $v$ (\cref{ch:g12:mechanical-waves})، مركزها النقطة التي
\emph{صدرت} منها --- والوسط لا يعرف أن المنبع تحرك ولا يعنيه ذلك.

\begin{proposition}[منبع متحرك]\label{prop:g12:doppler-effect:moving-source}
منبع يصدر قممًا \omterm{def:g10:signals-and-waves:period}{بدور} $T$ (\omterm{def:g10:signals-and-waves:frequency}{وتواتر} $f = 1/T$) وهو يتحرك في خط
مستقيم بسرعة $v_s < v$ عبر وسط تنتقل فيه \omterm{def:g10:signals-and-waves:wave}{الموجة} بسرعة $v$. فيستقبل
راصد ساكن في الوسط
\[
\lambda' = (v \mp v_s)\,T, \qquad f' = f\,\frac{v}{v \mp v_s},
\]
والإشارتان العلويتان لمنبع مقترب، والسفليتان لمنبع مبتعد.
\end{proposition}

\begin{proof}
بين قمتين تتقدم الأولى $vT$ بينما يتقدم المنبع $v_s T$. فأمامه تصدر
القمة \omterm{def:g10:orders-of-magnitude:unit}{الثانية} $v_s T$ \emph{أقرب} إلى الأولى، أي $\lambda' = (v - v_s)T$؛
وخلفه يُضاف المقدار $v_s T$ نفسه. ثم تجتاح القمم راصدًا ساكنًا
بسرعة الوسط $v$، واحدة كل $\lambda'/v$ \omterm{def:g10:orders-of-magnitude:unit}{ثانية}، ومنه $f' = v/\lambda' = f v/(v \mp v_s)$.
\end{proof}

\begin{omfigure}
\begin{tikzpicture}[font=\small]
  \foreach \i/\r in {0/3.0, 1/2.25, 2/1.5, 3/0.75} {
    \draw[omDef] ({0.3*\i},0) circle ({\r});
    \fill[gray] ({0.3*\i},0) circle (1.1pt);
  }
  \fill (1.2,0) circle (2pt);
  \draw[-Stealth, omThm, thick] (1.2,0) -- (1.75,0)
    node[below right=-2pt] {$\vect{v}_s$};
  \draw[Stealth-Stealth, thick] (2.55,0) -- (3.0,0);
  \node[above] at (2.78,0.05) {$\lambda'$};
  \draw[Stealth-Stealth, thick] (-3.0,0) -- (-1.95,0)
    node[midway, above] {$\lambda'$};
  \node[gray] at (0.15,0.55) {\scriptsize المواضع السابقة};
\end{tikzpicture}

{\small جبهات \omterm{def:g10:signals-and-waves:wave}{موجة} منبع يتحرك يمينًا: كل دائرة مركزها النقطة التي
صدرت منها (النقاط الرمادية)، ومن ثم تتزاحم القمم أمامه، $\lambda' = (v - v_s)T$،
وتتمدد خلفه، $\lambda' = (v + v_s)T$.}
\end{omfigure}

\begin{example}[سيارة الإسعاف، بالأرقام]\label{ex:g12:doppler-effect:siren}
صفّارة تواترها $f = \qty{700}{Hz}$ تقترب بسرعة $v_s = \qty{25}{m/s}$؛ وخذ
$v = \qty{340}{m/s}$ في الهواء (\cref{ch:g12:sound-acoustics}). عند الاقتراب:
$f' = 700 \times 340/315 \approx \qty{756}{Hz}$؛ وعند الابتعاد:
$f' = 700 \times 340/365 \approx \qty{652}{Hz}$. فيُنزل المرور الحدّةَ بمقدار \qty{104}{Hz}
--- أي بنسبة $\approx 1.16$، نحو نصفَي نغمة ونصف، لا تخطئها أي أذن.
\end{example}

\begin{remark}[ما لا تغيّره الحركة]\label{rem:g12:doppler-effect:speed-unchanged}
تظل \omterm{def:g10:signals-and-waves:wave}{الموجة} تنتقل بسرعة $v$: فالوسط وحده يحدّد السرعة
(\cref{ch:g12:mechanical-waves}). وحركة المنبع تعصر \emph{أطوال \omterm{def:g10:signals-and-waves:wave}{الموجة}}، لا
السرعات أبدًا. وللعصر حدّ: فحين $v_s \to v$ تتراكم القمم أمامه بعضها على بعض
و$f' \to \infty$ --- أي ازدحام جبهات \omterm{def:g10:signals-and-waves:wave}{الموجة} الذي تجرّه طائرة أسرع من \omterm{def:g10:signals-and-waves:sound}{الصوت} موجةً صدمية.
\end{remark}

\section{راصد متحرك، واختصار الحركة البطيئة}

\begin{proposition}[راصد متحرك]\label{prop:g12:doppler-effect:moving-observer}
يتحرك راصد بسرعة $v_o$ مستقيمًا نحو منبع ساكن في الوسط (أو مبتعدًا عنه).
فيكون \omterm{def:g10:signals-and-waves:frequency}{التواتر} المستقبَل
\[ f' = f\left(1 \pm \frac{v_o}{v}\right) \quad (+\ \text{مقتربًا},\ -\ \text{مبتعدًا}). \]
\end{proposition}

\begin{proof}
نمط \omterm{def:g10:signals-and-waves:wave}{الموجة} نفسه لم يضطرب: فالقمم المتباعدة $\lambda = vT$ تنتقل بسرعة $v$. والراصد
الراكض نحوها يلاقي القمم بالسرعة النسبية $v + v_o$، واحدة كل
$\lambda/(v + v_o)$ \omterm{def:g10:orders-of-magnitude:unit}{ثانية}: ومنه $f' = (v + v_o)/\lambda = f(1 + v_o/v)$. وبالركض بعيدًا، $v - v_o$.
\end{proof}

\begin{definition}[السرعة الشعاعية]\label{def:g12:doppler-effect:radial}
\emph{السرعة الشعاعية}\index{سرعة شعاعية} $v_r$ لمنبع بالنسبة إلى
راصد هي مركّبة سرعتهما النسبية على المستقيم الواصل بينهما --- أي
معدل تقلص المسافة أو نموّها. وهذه المركّبة وحدها تزيح \omterm{def:g10:signals-and-waves:frequency}{التواتر}؛
أما الحركة \emph{عرضًا} على خط النظر فلا تعطي انزياحًا في هذه الرتبة.
\end{definition}

\begin{proposition}[الحركة البطيئة: صيغة واحدة للجميع]\label{prop:g12:doppler-effect:slow}
حين تكون \omterm{def:g12:doppler-effect:radial}{السرعة الشعاعية} صغيرة، $v_r \ll v$، لم يعد يهمّ من الذي يتحرك:
\[
\frac{\Delta f}{f} = \frac{f' - f}{f} \approx +\frac{v_r}{v} \ \text{(اقتراب)},
\qquad
\frac{\Delta f}{f} \approx -\frac{v_r}{v} \ \text{(ابتعاد)}.
\]
\end{proposition}

\begin{proof}
صيغتا الراصد هما بالضبط $1 \pm v_o/v$ أصلًا. أما في حالة المنبع، فاكتب $x = v_s/v$:
فتبيّن المتطابقة $\frac{1}{1 - x} = 1 + x + \frac{x^2}{1 - x}$ أن $f'/f$ يختلف عن
$1 + x$ بحدّ من رتبة $x^2$ --- أي، من أجل $x = 0.1$، تصحيحًا قدره $1\%$ على انزياح قدره $10\%$.
وفي الرتبة الأولى تنهار الصيغ الأربع كلها إلى $1 \pm v_r/v$.
\end{proof}

\begin{omfigure}
\begin{tikzpicture}
\begin{axis}[omaxis, width=8.8cm, height=4.8cm,
    xlabel={$t$ (\unit{s})}, ylabel={$f'$ (\unit{Hz})},
    xmin=-8, xmax=8, ymin=408, ymax=472,
    ytick={420,440,460}, xtick={-8,-4,0,4,8}]
  \addplot[dashed, omExo, domain=-8:8, samples=2] {460.3};
  \addplot[dashed, omExo, domain=-8:8, samples=2] {421.4};
  \addplot[dotted, omProp, thick, domain=-8:8, samples=2] {440};
  \addplot[omDef, thick, domain=-8:8, samples=120]
    {440*340/(340 + 15*(15*x)/sqrt((15*x)^2 + 64))};
\end{axis}
\end{tikzpicture}

{\small \omterm{def:g10:signals-and-waves:frequency}{التواتر} المسموع عند مرور صفّارة تواترها \qty{440}{Hz} بسرعة \qty{15}{m/s}، على بعد \qty{8}{m}
من اللاقط: هضبتان عند $fv/(v \mp v_s)$، \omterm{def:g10:signals-and-waves:frequency}{والتواتر} الحقيقي $f$ (المنقّط) يُعبَر
أساسًا عند أقرب اقتراب، حيث تكون الحركة عرضية خالصة.}
\end{omfigure}

\begin{method}[قراءة مرور]\label{met:g12:doppler-effect:pass}
يُظهر تسجيل منبع مارّ هضبتين: $f_1$ (اقتراب) و$f_2$
(ابتعاد). ومنه
\[
f = \frac{2 f_1 f_2}{f_1 + f_2}, \qquad
v_s = v\,\frac{f_1 - f_2}{f_1 + f_2}.
\]
وفعلًا $f_1 + f_2 = \frac{2fv^2}{v^2 - v_s^2}$ و$f_1 - f_2 = \frac{2fv\,v_s}{v^2 - v_s^2}$:
فاقسم لتحصل على $v_s$، واحسب $2f_1 f_2/(f_1 + f_2)$ لتحصل على $f$. فقراءتا \omterm{def:g10:signals-and-waves:frequency}{تواتر} تعطيان
\omterm{def:g12:sound-acoustics:pitch-loudness}{الحدّة} \omterm{def:g11:lenses-and-eye:images}{الحقيقية} والسرعة معًا --- بلا ساعة إيقاف ولا مسطرة.
\end{method}

\section{أصداء تقيس السرعة}

اقذف \omterm{def:g10:signals-and-waves:wave}{موجة} على هدف متحرك يضرب \omterm{def:g12:doppler-effect:doppler}{مفعول دوبلر} مرتين.

\begin{proposition}[الانعكاس يضاعف الانزياح]\label{prop:g12:doppler-effect:double}
تُرسَل \omterm{def:g10:signals-and-waves:wave}{موجة} تواترها $f$ نحو عاكس يقترب مواجهةً بسرعة $u \ll v$.
فيعود \omterm{rem:g10:signals-and-waves:honest}{الصدى} \omterm{def:g10:signals-and-waves:frequency}{بتواتر}
\[ f' = f\,\frac{v + u}{v - u}, \qquad \text{ومنه} \qquad \Delta f \approx \frac{2u}{v}\,f . \]
\end{proposition}

\begin{proof}
انزياحان على التسلسل. فالعاكس، بوصفه \emph{راصدًا} متحركًا، يستقبل $f(1 + u/v)$؛
وبإعادة إصدار ما يستقبله يصير الآن \emph{منبعًا} متحركًا، فيصل \omterm{rem:g10:signals-and-waves:honest}{الصدى} عند
$f' = f(1 + u/v)/(1 - u/v) \approx f(1 + 2u/v)$ من أجل $u \ll v$.
\end{proof}

\begin{example}[مسدس الرادار]\label{ex:g12:doppler-effect:radar}
يصدر رادار مروري موجات مكروية عند $f = \qty{24.15}{GHz}$ ($v = c = \qty{3.00e8}{m/s}$)
ويركّب \omterm{rem:g10:signals-and-waves:honest}{الصدى} على \omterm{def:g10:signals-and-waves:wave}{الموجة} الصادرة؛ فتُسمع الضربة (\cref{ch:g12:sound-acoustics})
مباشرةً $\Delta f = 2uf/c \approx \qty{161}{Hz}$ لكل \unit{m/s} من سرعة السيارة
--- أي تواترًا سمعيًّا يُقاس بلا عناء منحوتًا من حامل تواتره \qty{24}{GHz}. والعامل
$2$ ليس اختياريًّا: فنسيانه يجامل كل سائق بالنصف.
\end{example}

\begin{example}[الموجات فوق الصوتية بدوبلر]\label{ex:g12:doppler-effect:ultrasound}
يرسل مسبار طبي \omterm{def:g10:signals-and-waves:wave}{موجة} فوق صوتية تواترها $f = \qty{5.0}{MHz}$ إلى شريان
($v \approx \qty{1540}{m/s}$ في النسيج)؛ وتعكسها كريات الدم الحمراء. فيعيد الدم عند
$u = \qty{0.50}{m/s}$ تواترًا $\Delta f = 2 \times 0.50 \times \num{5.0e6}/1540 \approx
\qty{3.2}{kHz}$ --- أي ضربة \emph{مسموعة}: فيسمع طبيب القلب حرفيًّا الدم
يتسارع عند كل نبضة قلب، ثم يقيس $u$ من $\Delta f$.
\end{example}

\section{دوبلر في ضوء النجوم}

الضوء \omterm{def:g10:signals-and-waves:wave}{موجة} (\cref{ch:g12:light-as-wave})، وطيفه مخطَّط بخطوط حادّة عند
أطوال \omterm{def:g10:signals-and-waves:wave}{موجة} يثبّتها كل عنصر كيميائي (\cref{ch:g10:light-spectra}) --- أي مسطرة
معايَرة في المختبر مطبوعة على كل نجم.

\begin{definition}[الانزياح نحو الأحمر ونحو الأزرق]\label{def:g12:doppler-effect:redshift}
حين تظهر الخطوط الطيفية لمنبع عند أطوال \omterm{def:g10:signals-and-waves:wave}{موجة} \emph{أطول} منها في
المختبر، يكون \omterm{def:g12:sound-acoustics:timbre}{الطيف} \emph{منزاحًا نحو الأحمر}\index{انزياح نحو الأحمر}؛ وحين تظهر \emph{أقصر}،
يكون \emph{منزاحًا نحو الأزرق}\index{انزياح نحو الأزرق} --- أي منزاحًا نحو الطرف الأحمر أو الأزرق من
\omterm{def:g10:light-spectra:wavelength}{الطيف المرئي}.
\end{definition}

\begin{proposition}[انزياح دوبلر للضوء]\label{prop:g12:doppler-effect:light}
في حالة منبع مبتعد \omterm{def:g12:doppler-effect:radial}{بسرعة شعاعية} $v_r \ll c$، يُرصد كل طول \omterm{def:g10:signals-and-waves:wave}{موجة}
ممدودًا بمقدار
\[ \frac{\Delta \lambda}{\lambda} = \frac{\lambda' - \lambda}{\lambda} \approx \frac{v_r}{c} \quad \text{(انزياح نحو الأحمر)}, \]
ومضغوطًا بالمقدار نفسه، $\Delta\lambda/\lambda \approx -v_r/c$، في حالة منبع
مقترب (\omterm{def:g12:doppler-effect:redshift}{انزياح نحو الأزرق}).
\end{proposition}
\admitted

\begin{remark}[أين تقيم الصيغة الأمينة]\label{rem:g12:doppler-effect:honest}
الضوء لا يحتاج وسطًا، فلا يمكن نسخ اشتقاق \omterm{def:g10:signals-and-waves:sound}{الصوت} --- إذ لا وجود لعبارة ``ساكن
في الوسط''. وتأتي الصيغة الدقيقة من النسبية الخاصة، لاحقًا هذه السنة
(\cref{ch:g12:special-relativity})؛ والتصحيحات من رتبة $(v_r/c)^2$، وهي مهمَلة في
الطائرات والنجوم والمجرّات القريبة. وفيما يلي نستعمل صيغة الحركة البطيئة بضمير مرتاح.
\end{remark}

\begin{omfigure}
\begin{tikzpicture}[font=\small]
  \foreach \y/\xl/\c/\t in {2.4/1.8/black/laboratory,
      1.2/0.6/omDef/approaching star, 0/3.72/omThm/receding star} {
    \fill[gray!18] (0,\y) rectangle (6,\y+0.55);
    \draw[\c, very thick] (\xl,\y) -- (\xl,\y+0.55);
    \node[right] at (6.15,\y+0.28) {\t};
  }
  \draw[dashed, gray] (1.8,-0.15) -- (1.8,2.95);
  \draw[-Stealth] (-0.2,-0.55) -- (6.4,-0.55)
    node[below left] {$\lambda$ (\unit{nm})};
  \foreach \x/\l in {0/656.0, 3/656.5, 6/657.0}
    \draw (\x,-0.48) -- (\x,-0.62) node[below] {\l};
\end{tikzpicture}

{\small خط الهيدروجين نفسه في ثلاثة أطياف: عند \qty{656.30}{nm} في المختبر،
\omterm{def:g12:doppler-effect:redshift}{ومنزاحًا نحو الأزرق} في نجم مقترب، ونحو الأحمر في نجم مبتعد. ويعطي الانزياح، مقروءًا
في مقابل \omterm{def:g10:relative-motion:frame}{المرجع} المتقطع، السرعةَ الشعاعية.}
\end{omfigure}

\begin{example}[وزن إفلات نجم]\label{ex:g12:doppler-effect:vega}
في \omterm{def:g12:sound-acoustics:timbre}{طيف} النجم الساطع النسر الواقع، يُقاس خط الهيدروجين الذي \omterm{def:g10:light-spectra:wavelength}{طول موجته} المخبري
\qty{656.30}{nm} عند \qty{656.27}{nm}: فيكون $\Delta\lambda = \qty{-0.03}{nm}$، أي
انزياحًا نحو الأزرق، ومن ثم يقترب النسر الواقع بسرعة $v_r = c\,|\Delta\lambda|/\lambda \approx \qty{1.4e4}{m/s}$
--- أي \qty{14}{km/s}، مقروءةً من انزياح قدره جزء من عشرين ألفًا.
\end{example}

\begin{example}[النجوم المتمايلة: الثنائيات والكواكب الخارجية]\label{ex:g12:doppler-effect:wobble}
حين يدور نجمان أحدهما حول الآخر، تتأرجح خطوطهما دوريًّا زرقاءَ وحمراء --- فيكونان
\emph{ثنائيًّا طيفيًّا}\index{ثنائي طيفي}، وتُقاس مداراتهما دون تفريق الزوج قط. ويفعل كوكب
بنجمه الشيء نفسه، بخفوت: فكلاهما يدور حول مركز كتلتهما المشترك، ومن ثم تتذبذب
\omterm{def:g12:doppler-effect:radial}{السرعة الشعاعية} للنجم ببضع عشرات من \unit{m/s} --- أي بانزياح قدره $\Delta\lambda/\lambda \sim 10^{-7}$.
وكشف ذلك هو ما وُجد به أول كوكب حول نجم شبيه بالشمس؛ ومسألة نهاية الأسبوع
تعيد ذلك الاكتشاف، بالأرقام كلها.
\end{example}

\begin{remark}[انزياحات المجرّات نحو الأحمر]\label{rem:g12:doppler-effect:expansion}
أطياف المجرّات البعيدة كلها منزاحة نحو \emph{الأحمر}، وكلما ازدادت بعدًا ازداد
انزياحها: فمسافات الكون تتمدد. وفي المجرّات القريبة تعطي
$v_r = c\,\Delta\lambda/\lambda$ سرعة الابتعاد؛ أما في الأبعد فتتجاوز الانزياحات
ما تقدر صيغة الحركة البطيئة على معالجته بأمانة، وتُستأنف المحاسبة
من \cref{ch:g12:special-relativity} فصاعدًا، وفي المجلدات الجامعية.
\end{remark}

\section{تمارين}

\begin{exercise}[$\star$]\label{exo:g12:doppler-effect:1}
تصدر صفّارة سيارة إطفاء عند \qty{700}{Hz} وتسير بسرعة \qty{25}{m/s} ($v = \qty{340}{m/s}$).
احسب \omterm{def:g10:signals-and-waves:frequency}{التواتر} الذي يسمعه ماشٍ (أ) أمامها؛ (ب) خلفها. (ج) وماذا يسمع
السائق؟
\end{exercise}

\begin{exercise}[$\star$]\label{exo:g12:doppler-effect:2}
يزعم صديق: ``تهبط \omterm{def:g12:sound-acoustics:pitch-loudness}{الحدّة} كلما ابتعدت سيارة الإسعاف لأن \omterm{def:g10:signals-and-waves:sound}{الصوت} يضعف
مع المسافة.'' فكّك الأثرين المختلطين، وقل على أيّ شيء يتوقف كلٌّ منهما.
\end{exercise}

\begin{exercise}[$\star$]\label{exo:g12:doppler-effect:3}
يركب راكب دراجة بسرعة \qty{6.0}{m/s} مستقيمًا نحو صفّارة ساكنة تصدر عند \qty{700}{Hz}.
فأيّ \omterm{def:g10:signals-and-waves:frequency}{تواتر} يسمع؟ وإذا ركب مبتعدًا مستقيمًا؟ وأيّ صيغة تنطبق، ولماذا لا تنطبق
صيغة المنبع المتحرك؟
\end{exercise}

\begin{exercise}[$\star$]\label{exo:g12:doppler-effect:4}
يصدر بوق قطار عند \qty{400}{Hz} بينما يسير القطار بسرعة \qty{30}{m/s}. احسب
\omterm{def:g12:mechanical-waves:wavelength}{طول الموجة} أمام القطار، وخلفه، وفي حالة السكون. وأيّ راصد يستقبل أيًّا منها؟
\end{exercise}

\begin{exercise}[$\star$]\label{exo:g12:doppler-effect:5}
هاتف يعزف نغمة تواترها \qty{440}{Hz} يحمله عدّاء بسرعة \qty{5.0}{m/s}. استعمل
صيغة الحركة البطيئة لتقدير الانزياح الذي يسمعه من يقترب منه العدّاء. وقارنه
بنصف نغمة (نحو $6\%$ في \omterm{def:g10:signals-and-waves:frequency}{التواتر}): فهل يلاحظه موسيقي؟
\end{exercise}

\begin{exercise}[$\star\star$]\label{exo:g12:doppler-effect:6}
يطير خفاش بسرعة \qty{6.0}{m/s} مستقيمًا نحو جدار، ويصدر عند \qty{50.0}{kHz}. بيّن أن
\omterm{rem:g10:signals-and-waves:honest}{الصدى} الذي يسمعه يعود عند $f' = f(v + v_b)/(v - v_b)$، ثم احسب $f'$ والانزياح.
ولماذا تكون هذه صيغة مسدس \omterm{met:g10:signals-and-waves:echo}{الرادار} متنكّرة؟
\end{exercise}

\begin{exercise}[$\star\star$]\label{exo:g12:doppler-effect:7}
مسدس رادار عند \qty{24.15}{GHz} يقيس ضربة قدرها \qty{4.0}{kHz} من سيارة مقتربة.
جِد سرعة السيارة \omterm{def:g10:orders-of-magnitude:unit}{بوحدة} \unit{m/s} وبوحدة \unit{km/h}. وأيّ ضربة تعطيها سيارة عند
\qty{50}{km/h} بالضبط؟
\end{exercise}

\begin{exercise}[$\star\star$]\label{exo:g12:doppler-effect:8}
مسبار دوبلر \omterm{def:g10:signals-and-waves:ultrasound}{فوق صوتي} عند \qty{5.0}{MHz} مصوَّب على امتداد شريان ($v = \qty{1540}{m/s}$ في
النسيج) يسمع ضربة قدرها \qty{2.6}{kHz}. جِد سرعة الدم. ولماذا يجب أن يُمال المسبار
\emph{على امتداد} الجريان لا عموديًّا عليه؟
\end{exercise}

\begin{exercise}[$\star\star$]\label{exo:g12:doppler-effect:9}
خذ $f = \qty{1000}{Hz}$ و$v = \qty{340}{m/s}$ وسرعة اقتراب قدرها \qty{34}{m/s}.
احسب $f'$ بدقة حين (أ) يتحرك المنبع، (ب) يتحرك الراصد. وقارن كليهما
بتنبؤ الحركة البطيئة $f(1 + v_r/v)$ واشرح حجم الاختلاف.
\end{exercise}

\begin{exercise}[$\star\star$]\label{exo:g12:doppler-effect:10}
يُظهر تسجيل دراجة نارية صغيرة تمرّ بلاقط هضبتين عند \qty{465}{Hz} و
\qty{415}{Hz}. باستعمال \cref{met:g12:doppler-effect:pass}، جِد \omterm{def:g10:signals-and-waves:frequency}{التواتر} الحقيقي للبوق و
سرعة الدراجة \omterm{def:g10:orders-of-magnitude:unit}{بوحدة} \unit{km/h}. ولماذا \emph{ليس} \omterm{def:g10:signals-and-waves:frequency}{التواتر} الحقيقي متوسطهما؟
\end{exercise}

\begin{exercise}[$\star\star$]\label{exo:g12:doppler-effect:11}
في \omterm{def:g12:sound-acoustics:timbre}{طيف} نجم، يُرصد خط الهيدروجين الذي \omterm{def:g10:light-spectra:wavelength}{طول موجته} المخبري \qty{656.3}{nm}
عند \qty{656.5}{nm}. أانزياح نحو الأحمر أم نحو الأزرق؟ احسب \omterm{def:g12:doppler-effect:radial}{السرعة الشعاعية} للنجم وجهتها.
\end{exercise}

\begin{exercise}[$\star\star\star$]\label{exo:g12:doppler-effect:12}
يتمايل نجم بسرعة \qty{55}{m/s} بسبب كوكب غير مرئي. احسب \omterm{def:g10:signals-and-waves:amplitude}{سعة}
تأرجح \omterm{def:g12:mechanical-waves:wavelength}{طول الموجة} لخط عند \qty{500}{nm}، طولًا وكسرًا من $\lambda$. \omterm{def:g10:orders-of-magnitude:oom}{ورتبة
القدر}: لماذا وجب أن ينتظر صيد الكواكب مطيافات مستقرة إلى جزء من
$10^{7}$؟
\end{exercise}

\begin{exercise}[$\star\star\star$]\label{exo:g12:doppler-effect:13}
يُوجَد كل خط في \omterm{def:g12:sound-acoustics:timbre}{طيف} مجرّة ممدودًا بمقدار $2.4\%$. احسب سرعة ابتعادها.
ويجد الفلكيون أن هذه السرعات تنمو تناسبًا مع المسافة، في كل اتجاه من
السماء: فأيّ صورة يوحي بها ذلك؟ وهل تبقى صيغة الحركة البطيئة جديرة بالثقة هنا؟
\end{exercise}

\begin{exercise}[$\star\star\star$]\label{exo:g12:doppler-effect:14}
تدور الشمس: فتقترب حافة قرصها منا بنحو \qty{2.0}{km/s} بينما تبتعد
الأخرى بالسرعة نفسها. فماذا يفعل ذلك بالخط \qty{656.3}{nm} في الضوء الآتي من القرص
كله --- أيزيحه أم يعرّضه؟ احسب الأثر \omterm{def:g10:orders-of-magnitude:unit}{بوحدة} \unit{nm}.
\end{exercise}

\begin{exercise}[$\star\star\star$]\label{exo:g12:doppler-effect:15}
سيارة شرطة بسرعة \qty{40}{m/s}، وصفّارتها عند \qty{700}{Hz}، تطارد شاحنة تسير بسرعة \qty{30}{m/s}
في الاتجاه نفسه. (أ) برّر أن $f' = f(v - v_o)/(v - v_s)$ بالنسبة إلى سائق الشاحنة.
(ب) احسب $f'$. (ج) بيّن أنه لو ساوت الشاحنة سرعة السيارة لتلاشى الانزياح،
وقل لماذا يجب أن يكون ذلك كذلك.
\end{exercise}

\section{مسألة: صيد كوكب خارج المجموعة الشمسية}

\begin{problem}[{مسألة نهاية الأسبوع --- صيد كوكب خارجي: بوق قطار يعاير
الأذن، ومسدس رادار يتمرّن على الانزياح المزدوج، وبحلول ليلة الأحد يكون طيف نجم،
متمايلًا بجزء من خمسة ملايين، قد وزن كوكبًا لم يره أحد قط}]
\label{pb:g12:doppler-effect:1}
يقضي نادٍ فلكي نهاية أسبوع على فكرة واحدة --- انزياح دوبلر --- متمرّنًا عليها على الأرض،
ثم مصوّبًا إياها نحو نجم شبيه بالشمس. المعطيات: $v = \qty{340}{m/s}$ \omterm{prop:g12:sound-acoustics:speed}{لسرعة الصوت} في الهواء،
و$c = \qty{3.00e8}{m/s}$، و$G = \qty{6.67e-11}{N.m^2/kg^2}$؛ وكتلة النجم
$M = \qty{2.0e30}{kg}$. ونقبل، من ميكانيك المدارات الدائرية
(\cref{ch:g12:satellites-kepler}): أن كوكبًا كتلته $m$ على \omterm{def:g10:universal-gravitation:satellite}{مدار} دائري نصف قطره $r$
له \omterm{def:g10:signals-and-waves:period}{دور} وسرعة يعطيهما $r^3 = \frac{G M P^2}{4\pi^2}$ و$V = \frac{2\pi r}{P}$،
وأن النجم، وهو يدور حول \omterm{def:g11:forces-and-motion:com}{مركز الكتلة} المشترك، يتمايل بسرعة $K = \frac{m}{M}\,V$.

\medskip
\noindent\textbf{الجزء I --- الجمعة: معبر السكة.}
يمرّ قطار بسرعة ثابتة وبوقه يعمل؛ فيسجّل هاتف $f_1 = \qty{471}{Hz}$ وهو
يقترب، و$f_2 = \qty{411}{Hz}$ بعد مروره.
\begin{enumerate}
  \item اشرح، بصورة جبهات \omterm{def:g10:signals-and-waves:wave}{الموجة} المتزاحمة، لماذا $f_1 > f_2$ مع أن البوق لا يتغير قط.
  \item اكتب المعادلتين اللتين تربطان $f_1$ و$f_2$ \omterm{def:g10:signals-and-waves:frequency}{بالتواتر} الحقيقي $f$ وسرعة القطار $u$.
  \item بيّن أن $u = v\,\dfrac{f_1 - f_2}{f_1 + f_2}$ واحسبها \omterm{def:g10:orders-of-magnitude:unit}{بوحدة} \unit{m/s} وبوحدة \unit{km/h}.
  \item بيّن أن $f = \dfrac{2 f_1 f_2}{f_1 + f_2}$ واحسبه. ولماذا يكون $f$
        \emph{دون} متوسط $f_1$ و$f_2$ بقليل؟
  \item احسب أطوال \omterm{def:g10:signals-and-waves:wave}{الموجة} أمام القطار وخلفه، \omterm{def:g12:mechanical-waves:wavelength}{وطول الموجة} في حالة السكون.
\end{enumerate}

\medskip
\noindent\textbf{الجزء II --- السبت: مسدس \omterm{met:g10:signals-and-waves:echo}{الرادار}.}
تعرض \omterm{def:g10:signals-and-waves:period}{دورية} مرور رادارها: $f = \qty{24.125}{GHz}$، ويقرأ المسدس
الضربة بين \omterm{def:g10:signals-and-waves:wave}{الموجة} الصادرة \omterm{rem:g10:signals-and-waves:honest}{والصدى}.
\begin{enumerate}[resume]
  \item اشرح لماذا يُزاح \omterm{rem:g10:signals-and-waves:honest}{صدى} سيارة متحركة \emph{مرتين} --- وسمِّ الوظيفة التي تؤديها
        السيارة في كل انزياح.
  \item بيّن أن الضربة في حالة سيارة سرعتها $u \ll c$ هي $\Delta f \approx \dfrac{2u}{c}\,f$.
  \item وتعيد سيارة $\Delta f = \qty{3.55}{kHz}$: احسب $u$.
  \item حوّل إلى \unit{km/h} وقارن بحدّ \qty{80}{km/h} المعلَن هناك.
  \item وكم هرتزًا من الضربة يقابل \qty{1}{km/h}؟ علّق: ما الذي يجعل انزياحًا نسبيًّا
        ضئيلًا كهذا في \omterm{def:g10:signals-and-waves:wave}{موجة} تواترها \qty{24}{GHz} سهل القياس؟
\end{enumerate}

\medskip
\noindent\textbf{الجزء III --- الأحد: النجم المتمايل.}
ينزّل النادي \omterm{def:g12:doppler-effect:radial}{السرعة الشعاعية} المقيسة لنجم شبيه بالشمس، ليلةً بعد ليلة،
مرسومة أدناه.

\begin{omfigure}
\begin{tikzpicture}
\begin{axis}[omaxis, width=8.8cm, height=4.6cm,
    xlabel={$t$ (بالأيام)}, ylabel={$v_r$ (\unit{m/s})},
    xmin=0, xmax=8.4, ymin=-72, ymax=72,
    xtick={0,2.1,4.2,6.3,8.4},
    xticklabels={$0$,$2.1$,$4.2$,$6.3$,$8.4$}, ytick={-55,0,55}]
  \addplot[dashed, omExo, domain=0:8.4, samples=2] {55};
  \addplot[dashed, omExo, domain=0:8.4, samples=2] {-55};
  \addplot[omDef, thick, domain=0:8.4, samples=140] {55*sin(360*x/4.2)};
\end{axis}
\end{tikzpicture}

{\small \omterm{def:g12:doppler-effect:radial}{السرعة الشعاعية} للنجم على مدى اثنتي عشرة ليلة: جيب نظيف --- وهو
بصمة جسم على \omterm{def:g10:universal-gravitation:satellite}{مدار} دائري، منظور إليه من حافته.}
\end{omfigure}

\begin{enumerate}[resume]
  \item ما الذي يُقاس فعلًا، ليلةً بعد ليلة، في \emph{\omterm{def:g12:sound-acoustics:timbre}{طيف}} النجم؟ اذكر
        العلاقة المستعملة لتحويله إلى $v_r$.
  \item اقرأ \omterm{def:g10:signals-and-waves:amplitude}{سعة} $K$ \omterm{def:g10:signals-and-waves:period}{ودور} $P$ التمايل من المنحني.
  \item احسب التأرجح المقابل $\Delta\lambda$ لخط عند \qty{550}{nm}، و
        $\Delta\lambda/\lambda$. وقارن بالجزء I: كم يكون هذا القياس أصعب
        من قياس القطار؟
  \item ولماذا يجعل كوكب غير مرئي النجمَ يتمايل أصلًا؟ وعلامَ يدلّ الشكل
        الجيبي بشأن \omterm{def:g10:universal-gravitation:satellite}{المدار}؟
  \item ولماذا لا يقيس هذا الأسلوب إلا الجزء \emph{الشعاعي} من حركة النجم، وما
        الذي كان سيُرى لو كان \omterm{def:g10:universal-gravitation:satellite}{المدار} مواجهًا لنا؟
\end{enumerate}

\medskip
\noindent\textbf{الجزء IV --- ليلة الأحد: \omterm{def:g10:universal-gravitation:weight}{وزن} ما لا يُرى.}
خذ \omterm{def:g10:universal-gravitation:satellite}{المدار} دائريًّا ومنظورًا إليه من حافته، بحيث تكون $K$ السرعة المدارية الكاملة للنجم.
\begin{enumerate}[resume]
  \item انطلاقًا من $P$ والعلاقة المقبولة $r^3 = G M P^2/(4\pi^2)$، احسب نصف قطر
        \omterm{def:g10:universal-gravitation:satellite}{مدار} الكوكب $r$. وقارنه بالمسافة بين الأرض والشمس، \qty{1.5e11}{m}.
  \item احسب السرعة المدارية للكوكب $V = 2\pi r/P$.
  \item وانطلاقًا من $K = (m/M)\,V$، احسب كتلة الكوكب $m$.
  \item قارن $m$ بالمشتري ($\qty{1.9e27}{kg}$) وبالأرض ($\qty{6.0e24}{kg}$)، و
        $r$ بالمسافة بين الشمس وعطارد ($\qty{5.8e10}{m}$). فأيّ نوع من العوالم هذا؟
  \item ولو كان \omterm{def:g10:universal-gravitation:satellite}{المدار} مائلًا لا منظورًا إليه من حافته، فهل كانت الكتلة \omterm{def:g11:lenses-and-eye:images}{الحقيقية} أكبر أم
        أصغر من قيمتك؟ اختم في جملة واحدة: ما الذي سلّمه للتوّ انزياح دوري قدره جزء
        من خمسة ملايين؟
\end{enumerate}
\end{problem}
